\chapter{Homologia przestrzeni}
\label{ch:homologia-przestrzeni}

W Rozdziale~\ref{ch:algebra-homologiczna} homologia była
niezmiennikiem \emph{kompleksu łańcuchowego}.
Teraz przypiszemy taki kompleks przestrzeni. Najpierw
użyjemy sympleksów z Rozdziału~\ref{ch:topologia-algebraiczna-podstawy};
potem zastąpimy je wszystkimi ciągłymi obrazami sympleksów,
aby móc badać przestrzenie bez wybranej triangulacji. Na końcu
zobaczymy, dlaczego do obliczeń zwykle wystarczy
kilka komórek CW.

Intuicja jest jedna: cykl to suma zorientowanych kawałków
bez brzegu, a brzeg to cykl, który ogranicza kawałek
o wymiar wyższy. Iloraz
\[
 H_n(X;\Z)=
 \frac{\{\text{$n$-cykle w $X$}\}}
 {\{\text{brzegi $(n+1)$-łańcuchów w $X$}\}}
\]
odróżnia na przykład pętlę będącą brzegiem na tarczy od pętli
obiegającej dziurę. Pracujemy nad $\Z$; zmianę
współczynników omawia
Twierdzenie~\ref{thm:homologia-char-zero},
a konkretny rachunek znajduje się poniżej.

\section{Łańcuchy symplicjalne: orientacja i brzeg}
\label{sec:homologia-symplicjalna}

\begin{definicja}[Kompleks łańcuchów symplicjalnych]
\index{kompleks lancuchow symplicjalnych@Kompleks łańcuchów symplicjalnych}
\label{def:homologia-symplicjalna}
Niech $K$ będzie skończonym kompleksem symplicjalnym
z Definicji~\ref{def:kompleks-symplicjalny}.
Ustalmy porządek jego wierzchołków. Grupa
$C_n^\Delta(K)$ jest wolną grupą abelową z bazą
wszystkich $n$-sympleksów zapisanych w porządku
$[v_0,\ldots,v_n]$. Jeśli wierzchołki wypiszemy
w innej kolejności, przyjmujemy zasadę
\[
 [v_{\pi(0)},\ldots,v_{\pi(n)}]
 =\operatorname{sgn}(\pi)[v_0,\ldots,v_n].
\]
Zatem odwrócenie orientacji zmienia znak, a dwa
różne sympleksy nadal dają niezależne generatory.
Dla $n\geq1$ określamy
\begin{equation}\label{eq:brzeg-symplicjalny}
 \partial_n[v_0,\ldots,v_n]
 =\sum_{i=0}^n(-1)^i
 [v_0,\ldots,\widehat v_i,\ldots,v_n],
 \qquad \partial_0=0.
\end{equation}
Rozszerzamy $\partial$ liniowo na sumy
o całkowitych współczynnikach. Dla $n<0$ przyjmujemy $C_n^\Delta(K)=0$.
\end{definicja}

Znak ściany nie jest ozdobą: kierunek odziedziczony
przez brzeg wypełnionego trójkąta musi być spójny.
Dla $e=[v_0,v_1]$ i $t=[v_0,v_1,v_2]$ dostajemy
\[
 \partial e=[v_1]-[v_0],\qquad
 \partial t=[v_1,v_2]-[v_0,v_2]+[v_0,v_1].
\]
Drugą równość można zapisać
$\partial t=[v_0,v_1]+[v_1,v_2]+[v_2,v_0]$.
Jest to obieg brzegu trójkąta.

\begin{figure}[!htbp]
\centering
\begin{tikzpicture}[>=Stealth,line cap=round,line join=round]
 \coordinate (a) at (0,0); \coordinate (b) at (3.3,0);
 \coordinate (c) at (1.35,2.55);
 \fill[manifoldblue!13] (a)--(b)--(c)--cycle;
 \draw[manifoldblue!65!black,line width=1.6pt]
  (a)--(b)--(c)--cycle;
 \draw[->,manifoldred,line width=1.5pt]
  (.7,-.13)--(2.5,-.13);
 \draw[->,manifoldred,line width=1.5pt]
  (3.1,.38)--(1.8,2.07);
 \draw[->,manifoldred,line width=1.5pt]
  (1.04,2.1)--(.13,.39);
 \foreach \p in {a,b,c}{\fill[manifoldgold!85!black] (\p) circle(2.8pt);}
 \node[below left] at (a) {$v_0$};
 \node[below right] at (b) {$v_1$};
 \node[above] at (c) {$v_2$};
 \node[manifoldblue!75!black] at (1.4,.84) {$t$};
 \node[align=left,anchor=west,text width=5.5cm] at (4.1,1.38)
  {$\partial t=e_{01}+e_{12}+e_{20}$\\[3pt]
   zewnętrzna pętla jest brzegiem, gdy wnętrze $t$
   należy do przestrzeni};
\end{tikzpicture}
\caption{Zorientowany brzeg trójkąta. Usunięcie jego
wnętrza pozostawia cykl, który w samym okręgu
nie jest już brzegiem żadnego $2$-łańcucha.}
\label{fig:homologia-trojkat}
\end{figure}

\begin{twierdzenie}[Brzeg brzegu znika]
\label{thm:brzeg-brzegu}
Odwzorowania~\eqref{eq:brzeg-symplicjalny}
spełniają $\partial_{n-1}\partial_n=0$.
Wobec tego
\[
 H_n^\Delta(K):=
 \ker(\partial_n:C_n^\Delta(K)\to C_{n-1}^\Delta(K))
 \big/
 \operatorname{im}(\partial_{n+1})
\]
jest dobrze określoną grupą. Licznik składa
się z \emph{cykli}, mianownik z \emph{brzegów}.
\end{twierdzenie}
\begin{proof}
Dla $n=1$ wynika to z $\partial_0=0$.
Niech $n\geq2$ i weźmy generator
$[v_0,\ldots,v_n]$. Każda jego ściana
o wymiarze $n-2$ powstaje po usunięciu
pewnych dwóch wierzchołków $v_i,v_j$,
$i<j$, dokładnie na dwa sposoby.
Usunięcie najpierw $v_i$, a następnie
$v_j$ (ma on już pozycję $j-1$) daje
współczynnik $(-1)^i(-1)^{j-1}$.
W przeciwnej kolejności współczynnik
wynosi $(-1)^j(-1)^i$.
Ich suma to
$(-1)^{i+j-1}+(-1)^{i+j}=0$.
Dotyczy to każdej ściany z osobna,
więc $\partial^2$ znika na bazie,
a przez liniowość na całej grupie.
Stąd $\operatorname{im}\partial_{n+1}
\subseteq\ker\partial_n$, co pozwala
utworzyć podany iloraz.
\end{proof}

\begin{przyklad}[Odcinek, okrąg, tarcza]
\label{prz:homologia-odcinek-okrag}
Odcinek ma dwa wierzchołki $v_0,v_1$
i jedną krawędź $e$. Odwzorowanie
$\partial_1:\Z e\to\Z v_0\oplus\Z v_1$
przypisuje $me\mapsto m(v_1-v_0)$.
Ma zerowe jądro, więc $H_1^\Delta=0$,
a w ilorazie $H_0^\Delta$ obie klasy
wierzchołków są równe. Otrzymujemy $\Z$.

Teraz weźmy obwód trójkąta, bez wnętrza,
i ustawmy kierunki $e_{01},e_{12},e_{20}$.
Każdy $1$-łańcuch
$ae_{01}+be_{12}+ce_{20}$ jest cyklem
dokładnie wtedy, gdy $a=b=c$
(współczynnik przy każdym wierzchołku
w jej brzegu musi być zerem).
Nie ma tu $2$-sympleksów, a więc
$H_1^\Delta\cong\Z$ z generatorem
$e_{01}+e_{12}+e_{20}$.
Jeśli wstawimy wnętrze $t$, ten sam
cykl staje się $\partial t$ i
$H_1^\Delta$ znika. Rysunek~\ref{fig:homologia-trojkat}
pokazuje dokładnie tę zmianę.
\end{przyklad}

W $\Delta$-kompleksie z
Definicji~\ref{def:delta-kompleks} postępujemy
podobnie: każdy \emph{sympleks charakterystyczny}
jest osobnym generatorem, nawet jeżeli jego
różne ściany mają ten sam obraz. Brzeg
jest sumą obrazów ścian z ich znakami;
po identyfikacji mogą się one wzajemnie
skasować. Tak właśnie sklejony odcinek
z jednym wierzchołkiem ma $\partial e=v-v=0$.

\begin{stwierdzenie}[Składowe i $H_0$]
\label{prop:h0-symplicjalne}
Jeżeli $K$ ma $c$ składowych łukowych,
to $H_0^\Delta(K)\cong\Z^c$.
\end{stwierdzenie}
\begin{proof}
Ponieważ $\partial_0=0$, grupa $H_0$
to wolna grupa na wierzchołkach
modulo relacje
$v_j-v_i=\partial[v_i,v_j]$
dla krawędzi. Idąc po drodze krawędziowej,
utożsamiamy klasy wszystkich wierzchołków
jej składowej. Dwa wierzchołki tej samej
składowej $|K|$ można połączyć taką drogą:
wszystkie wierzchołki jednego sympleksu należą do tej samej
składowej grafu krawędzi. Sympleksy przypisane jednej takiej
składowej tworzą podkompleks. Jego realizacja i jej dopełnienie
są domknięte jako skończone sumy domkniętych sympleksów,
więc są też otwarte. Droga nie może przejść między tymi zbiorami.
Każdy punkt sympleksu można natomiast połączyć odcinkiem z wierzchołkiem.
Składowe grafu odpowiadają zatem dokładnie składowym łukowym $|K|$. Wreszcie żadna
relacja nie łączy różnych składowych,
bo obie końcówki każdej krawędzi leżą
w jednej. W każdej składowej pozostaje
więc jeden wolny generator.
\end{proof}

\section{Łańcuchy singularne i niezmienniczość homotopijna}
\label{sec:homologia-singularna}

Triangulacja bywa trudna do znalezienia. W konstrukcji
singularnej sam wybór \emph{ciągłego} odwzorowania
standardowego sympleksu do przestrzeni tworzy
generator; nie żądamy różnowartościowości ani
ładnego kształtu jego obrazu.

\begin{definicja}[Homologia singularna]
\index{homologia singularna@Homologia singularna}
\label{def:homologia-singularna}
\emph{Singularny $n$-sympleks} w $X$ to dowolne
ciągłe odwzorowanie $\sigma:\Delta^n\to X$.
Niech $C_n(X)$ będzie wolną grupą abelową
na wszystkich takich odwzorowaniach; dla $n<0$ przyjmujemy $C_n(X)=0$.
W szczególności
dwa różne odwzorowania o tym samym obrazie są
różnymi generatorami. Niech
$\iota_i:\Delta^{n-1}\to\Delta^n$
będzie afinicznym włączeniem ściany
naprzeciw $i$-tego wierzchołka. Kładziemy
\begin{equation}\label{eq:singular-boundary}
 \partial\sigma=\sum_{i=0}^n(-1)^i
 \sigma\circ\iota_i\quad(n\geq1),
 \qquad \partial:C_0(X)\to0,
 \qquad H_n(X):=H_n(C_\bullet(X),\partial).
\end{equation}
\end{definicja}

Obliczenie $\partial^2=0$ jest słowo w słowo
rachunkiem dwóch usuwanych wierzchołków
z dowodu Twierdzenia~\ref{thm:brzeg-brzegu}.
Teraz pojawia się nawet singularny sympleks
stały: nie odrzucamy go na poziomie
łańcuchów, gdyż później jego wkład może
być brzegiem.

\begin{figure}[!htbp]
\centering
\begin{tikzpicture}[>=Stealth,line cap=round,line join=round]
 \coordinate (a) at (0,0); \coordinate (b) at (2.45,0);
 \coordinate (c) at (.9,1.9);
 \fill[manifoldblue!12] (a)--(b)--(c)--cycle;
 \draw[manifoldblue!70!black,thick] (a)--(b)--(c)--cycle;
 \node[below left] at (a) {$v_0$};
 \node[below right] at (b) {$v_1$};
 \node[above] at (c) {$v_2$};
 \node at (1.05,.55) {$\Delta^2$};
 \draw[->,manifoldred,very thick] (2.95,.95)--(4.28,.95)
  node[midway,above] {$\sigma$};
 \fill[manifoldgreen!12]
  (4.95,.12) .. controls (5.7,-.65) and (6.6,.24) ..
  (7.75,-.08) .. controls (8.8,.3) and (8.3,1.4) ..
  (7.5,1.55) .. controls (6.5,2.2) and (5.2,1.66) ..
  (4.95,.12);
 \draw[manifoldgreen!70!black,thick]
  (4.95,.12) .. controls (5.7,-.65) and (6.6,.24) ..
  (7.75,-.08) .. controls (8.8,.3) and (8.3,1.4) ..
  (7.5,1.55) .. controls (6.5,2.2) and (5.2,1.66) ..
  (4.95,.12);
 \fill[manifoldblue!25,opacity=.7]
  (5.6,.45) .. controls (6.3,.1) and (7.3,.2) .. (7.9,.65)
  .. controls (8.05,1.1) and (7.3,1.6) .. (6.8,1.5)
  .. controls (6.1,1.7) and (5.8,1.05) .. (5.6,.45);
 \draw[manifoldred,thick,->]
  (5.6,.45) .. controls (6.3,.1) and (7.3,.2) .. (7.9,.65);
 \draw[manifoldred,thick,->]
  (7.9,.65) .. controls (8.05,1.1) and (7.3,1.6) .. (6.8,1.5);
 \draw[manifoldred,thick,->]
  (6.8,1.5) .. controls (6.1,1.7) and (5.8,1.05) .. (5.6,.45);
 \foreach \x/\y in {5.6/.45,7.9/.65,6.8/1.5}
  {\fill[manifoldred] (\x,\y) circle(1.8pt);}
 \node[font=\scriptsize,below left] at (5.6,.45) {$\sigma(v_0)$};
 \node[font=\scriptsize,right] at (7.9,.65) {$\sigma(v_1)$};
 \node[font=\scriptsize,above] at (6.8,1.5) {$\sigma(v_2)$};
 \node[manifoldgreen!70!black] at (8.05,1.87) {$X$};
 \node[manifoldblue!70!black] at (6.7,.8) {$\sigma(\Delta^2)$};
\end{tikzpicture}
\caption{Singularny sympleks jest odwzorowaniem do $X$, niekoniecznie
płaskim trójkątem. Jego brzeg liczymy
z ograniczeń odwzorowania $\sigma$ do trzech ścian.}
\label{fig:singular-simplex}
\end{figure}

\begin{stwierdzenie}[Ciągłe odwzorowania działają na homologii]
\label{prop:funktor-singularny}
Jeśli $f:X\to Y$ jest ciągłe, to
$f_\#:C_n(X)\to C_n(Y)$,
$f_\#(\sigma)=f\circ\sigma$,
jest odwzorowaniem łańcuchowym. Wyznacza homomorfizm
$f_*:H_n(X)\to H_n(Y)$, przy czym
$(g\circ f)_*=g_*\circ f_*$ oraz
$(\id_X)_*=\id_{H_n(X)}$.
\end{stwierdzenie}
\begin{proof}
Na pojedynczym generatorze
\[
 \partial(f\circ\sigma)
 =\sum_i(-1)^i f\circ\sigma\circ\iota_i
 =f_\#(\partial\sigma).
\]
Równość na sumach wynika z liniowości.
Zatem obraz cyklu jest cyklem, a obraz
brzegu brzegiem, więc $f_*$ jest
dobrze określony na ilorazie.
Równości dotyczące składania sprawdzamy
już na $\sigma$:
$g_\#f_\#(\sigma)=g\circ f\circ\sigma$.
\end{proof}

\begin{przyklad}[Punkt]
\label{prz:homologia-punktu}
Dla $X=\{*\}$ w każdym stopniu jest
dokładnie jeden singularny sympleks
stały $c_n$, więc $C_n(X)\cong\Z$.
Jego brzeg to
$\partial c_n=(\sum_{i=0}^n(-1)^i)c_{n-1}$:
dla $n$ nieparzystego współczynnik
jest $0$, dla dodatniego parzystego $1$.
Kompleks ma więc postać
\[
 \cdots\xrightarrow{\;0\;}\Z
 \xrightarrow{\;1\;}\Z
 \xrightarrow{\;0\;}\Z
 \xrightarrow{\;1\;}\Z
 \xrightarrow{\;0\;}\Z\longrightarrow0,
\]
gdzie prawy skraj to $C_0$; precyzyjnie
$\partial_1=0$, $\partial_2=1$.
Stąd $H_0(*)=\Z$ i $H_n(*)=0$ dla $n>0$.
\end{przyklad}

\begin{stwierdzenie}[Składowe łukowe]
\label{prop:h0-singularne}
Grupa $H_0(X)$ jest wolna na składowych
łukowych $X$. W szczególności
dla przestrzeni niepustej łukowo spójnej
$H_0(X)\cong\Z$.
\end{stwierdzenie}
\begin{proof}
Każdy punkt $x$ jest singularnym
$0$-sympleksem, więc $C_0(X)$
jest wolną grupą na punktach.
Dla drogi $\gamma:[0,1]=\Delta^1\to X$
mamy $\partial\gamma=[\gamma(1)]-
[\gamma(0)]$. Punkty w jednej składowej
łukowej mają zatem tę samą klasę.
Z drugiej strony każdy składnik
brzegu dowolnego $1$-łańcucha
jest różnicą punktów połączonych drogą:
nie może utożsamić dwóch różnych
składowych. Sumowanie współczynników
punktów osobno w każdej składowej
daje odwrotny homomorfizm do
wolnej grupy na składowych.
\end{proof}

\begin{twierdzenie}[Niezmienniczość homotopijna]
\label{thm:homotopia-singularna}
Jeśli $f,g:X\to Y$ są homotopijne, to
$f_*=g_*:H_n(X)\to H_n(Y)$ w każdym stopniu.
Jeśli $X$ i $Y$ są homotopijnie równoważne,
ich grupy homologii są izomorficzne.
\end{twierdzenie}
\begin{proof}
Niech $F:X\times[0,1]\to Y$ będzie
homotopią, $F(x,0)=f(x)$,
$F(x,1)=g(x)$.
Graniastosłup $\Delta^n\times[0,1]$
dzielimy na $n+1$ zorientowanych
$(n+1)$-sympleksów. Dla $i=0,\ldots,n$
ich uporządkowane wierzchołki to
\[
 (e_0,0),\ldots,(e_i,0),
 (e_i,1),\ldots,(e_n,1).
\]
Niech $T_i:\Delta^{n+1}\to
\Delta^n\times[0,1]$ będzie odwzorowaniem
afinicznym posyłającym kolejne wierzchołki
właśnie na tę listę. Określmy
\[
 P_n(\sigma)=\sum_{i=0}^n(-1)^i\,
 F\circ(\sigma\times\id_{[0,1]})\circ T_i
 \in C_{n+1}(Y).
\]
Rozpiszmy brzeg sumy. Wewnętrzne
$n$-ściany, na których spotykają się
sąsiednie części graniastosłupa,
pojawiają się z przeciwnymi znakami.
Ściany boczne nad każdą ścianą
$\Delta^n$ tworzą dokładnie
$-P_{n-1}(\partial\sigma)$:
jeśli $j<i$, usuwamy dolny wierzchołek $(e_j,0)$
z $T_i$, z łącznym znakiem $(-1)^{i+j}$.
Na pryzmacie ściany naprzeciw $e_j$ punkt podziału ma indeks $i-1$,
więc znak w $-P_{n-1}\partial$ wynosi
$-(-1)^j(-1)^{i-1}=(-1)^{i+j}$.
Jeśli $j>i$, usuwamy górny wierzchołek $(e_j,1)$,
na pozycji $j+1$; oba rachunki dają $(-1)^{i+j+1}$.
Wewnętrzna ściana między $T_{i-1}$ i $T_i$ powstaje
przez usunięcie odpowiednio $(e_{i-1},1)$ i $(e_i,0)$,
ze znakami $-1$ i $+1$.
Pozostałe dwie ściany to górna
$g\circ\sigma$ ze znakiem $+1$
oraz dolna $f\circ\sigma$ ze
znakiem $-1$. Zatem
\begin{equation}\label{eq:prism-identity}
 \partial P_n+P_{n-1}\partial
 =g_\#-f_\#.
\end{equation}
Przyjmujemy $P_{-1}=0$. Dla $n=0$ jest to znana równość
,,koniec drogi minus początek'';
powyższe parowanie ścian dowodzi
jej w każdym wymiarze. W świetle
Twierdzenia~\ref{thm:homotopia-homologia}
oba odwzorowania indukują to samo na homologii.
Jeśli $u:X\to Y$, $v:Y\to X$
są odwrotne z dokładnością do homotopii,
to $(vu)_*=\id$ i $(uv)_*=\id$,
więc $u_*$ i $v_*$ są odwrotne.
\end{proof}

W szczególności przestrzeń ściągalna
ma homologię punktu. Jeśli
$A\subseteq X$ jest retraktem
deformacyjnym, to $H_n(A)\cong H_n(X)$.

\section{Pary przestrzeni i homologia zredukowana}
\label{sec:homologia-wzgledna}

\begin{definicja}[Grupy względne i zredukowane]
\index{grupy wzgledne i zredukowane@Grupy względne i zredukowane}
\label{def:homologia-wzgledna}
Dla $A\subseteq X$ włączenie daje
podkompleks $C_\bullet(A)\subseteq C_\bullet(X)$.
Kładziemy
\[
 C_n(X,A)=C_n(X)/C_n(A),\qquad
 H_n(X,A)=H_n(C_\bullet(X,A)).
\]
\emph{Augmentacja} $\varepsilon:C_0(X)\to\Z$ sumuje
współczynniki punktów. Dołączamy stopień $C_{-1}=\Z$,
z różniczką $\partial_0=\varepsilon$; niższe grupy są zerowe.
Ponieważ $\varepsilon\partial_1=0$, otrzymujemy kompleks,
którego homologie oznaczamy $\widetilde H_n(X)$.
Dla niepustego $X$ mamy
$\widetilde H_0(X)=\ker(\varepsilon:H_0(X)\to\Z)$,
$\widetilde H_n(X)=H_n(X)$ dla $n>0$ i
$\widetilde H_n(X)=0$ dla $n<0$.
Dla $X=\varnothing$ wszystkie grupy łańcuchów w stopniach
nieujemnych są zerowe, więc $\widetilde H_{-1}(\varnothing)=\Z$
i pozostałe grupy zredukowane są zerowe.
\end{definicja}

Względny cykl może mieć brzeg w $A$:
to, co w $X$ jest otwartą powierzchnią,
po ,,pominięciu $A$'' staje się cyklem.
Na przykład zorientowany dysk reprezentuje klasę względną
względem swojego okręgu brzegowego.
Jeżeli $X$ ma skończoną liczbę $c>0$ składowych łukowych,
to $\widetilde H_0(X)\cong\Z^{c-1}$.
Ogólnie, po wyborze jednej składowej $P_0$, bazę jądra
augmentacji tworzą różnice $[P]-[P_0]$ dla $P\ne P_0$:
jest to suma prosta kopii $\Z$, także dla nieskończenie wielu składowych.

\begin{twierdzenie}[Długi ciąg pary]
\label{thm:les-pary}
Każda para $(X,A)$ daje naturalny
długi ciąg dokładny
\begin{equation}\label{eq:les-pary}
 \cdots\to H_n(A)\xrightarrow{i_*}H_n(X)
 \xrightarrow{j_*}H_n(X,A)
 \xrightarrow{\delta}H_{n-1}(A)
 \to H_{n-1}(X)\to\cdots.
\end{equation}
Jeśli $\partial c\in C_{n-1}(A)$,
to $\delta[c]=[\partial c]$.
\end{twierdzenie}
\begin{proof}
W każdej grupie łańcuchów mamy
krótki ciąg dokładny
$0\to C_n(A)\to C_n(X)\to C_n(X,A)\to0$.
Jest zgodny z brzegami, więc
Stwierdzenie~\ref{prop:ses-iloraz}
daje długi ciąg i wzór na $\delta$.
Naturalność pochodzi z przemienności
odwzorowań łańcuchów dla odwzorowań
par $(X,A)\to(Y,B)$.
\end{proof}

Jeśli $A\ne\varnothing$, możemy w tym ciągu zastąpić homologie
$A$ i $X$ zredukowanymi, pozostawiając zwykłą homologię względną.
Istotnie, w stopniu $-1$ włączenie kompleksów augmentowanych
jest identycznością $\Z\to\Z$, więc ich iloraz ma tam zero
i jest właśnie kompleksem $C_\bullet(X,A)$.

Homotopia par (w każdym czasie obraz $A$
leży w $B$) daje te same odwzorowania na $H_n(X,A)$:
pryzmat z dowodu Twierdzenia~\ref{thm:homotopia-singularna}
przenosi $C_\bullet(A)$ do $C_{\bullet+1}(B)$,
więc wzór~\eqref{eq:prism-identity}
schodzi na kompleksy ilorazowe.

\section{Podział barycentryczny, wycinanie i sklejanie}
\label{sec:wycinanie-homologia}

Homologia singularna ma nieskończenie wiele generatorów nawet
dla prostych przestrzeni. Podział dużego sympleksu na małe
części pozwala jednak sprowadzić rachunek do lokalnych
fragmentów. Z tego mechanizmu wynika wycinanie.

\begin{lemat}[Małe łańcuchy]
\label{lem:male-lancuchy}
Niech $\mathcal U=\{U_\alpha\}$ będzie otwartym
pokryciem $X$. Podkompleks
$C_n^{\mathcal U}(X)\subseteq C_n(X)$
generowany przez sympleksy, których obrazy
leżą w jednym $U_\alpha$, ma te same
grupy homologii co $C_\bullet(X)$.
Włączenie indukuje izomorfizmy w każdym stopniu.
\end{lemat}
\begin{proof}
Podamy konstrukcję potrzebną w obu kierunkach.
Podzielmy liniowo $\Delta^n$ na sympleksy
barycentryczne: wierzchołkami małego sympleksu
są barycentra łańcucha niepustych ścian
$F_0\subsetneq F_1\subsetneq\cdots\subsetneq F_n$,
zapisane w tej właśnie kolejności. Każdy taki łańcuch odpowiada
permutacji $\pi$ wierzchołków: $F_j=[e_{\pi(0)},\ldots,e_{\pi(j)}]$.
Niech $a_\pi:\Delta^n\to\Delta^n$ będzie afinicznym odwzorowaniem
na uporządkowany mały sympleks. Określamy
\[
 S_n(\sigma)=\sum_{\pi\in S_{n+1}}\operatorname{sgn}(\pi)
                 \,\sigma\circ a_\pi.
\]
Znak pochodzi z orientacji: dla permutacji identycznościowej
macierz kolejnych barycentrów ma dodatni wyznacznik w otoczce afinicznej,
a permutacja wierzchołków zmienia go o $\operatorname{sgn}(\pi)$.
Kolejność barycentrów jest ważna: w kompleksie singularnym
odwrócenie parametryzacji nie jest z definicji ujemnym generatorem.
Ściany wewnętrzne dwóch części
występują z przeciwnymi znakami,
a ich ściany na zewnątrz dają
podział odpowiedniej ściany $\sigma$.
Dokładniej, wewnętrzną ścianę uzyskujemy, pomijając barycentrum $F_j$
dla $j<n$. Dwie odpowiadające jej permutacje zamieniają miejscami
pozycje $j,j+1$, więc mają przeciwne znaki, a pozostałe
uporządkowane barycentra są identyczne. Dla $j=n$ otrzymujemy
sympleks podziału ściany naprzeciw $e_{\pi(n)}$.
Jeśli $\rho$ jest permutacją pozostałych wierzchołków w ich
odziedziczonej kolejności, to
$(-1)^n\operatorname{sgn}(\pi)=(-1)^{\pi(n)}\operatorname{sgn}(\rho)$.
To właśnie znak tej ściany w $S\partial$; stąd $\partial S=S\partial$.

Wyjaśnijmy też, dlaczego $S$ nie zmienia
klasy homologii. Na standardowym sympleksie
oznaczmy jego identycznościowe odwzorowanie
przez $u_n$ i indukcyjnie zbudujmy
łańcuchy $t_n$ tak, aby dla $n\geq1$
\[
 \partial t_n+\sum_{i=0}^n(-1)^i(\iota_i)_\#t_{n-1}
       =S(u_n)-u_n. \tag{*}
\]
W stopniu $0$ obie strony są zerowe,
więc $t_0=0$. Jeśli mamy już konstrukcję
na wszystkich ścianach, to oznaczmy ich sumę
$b_n=\sum_i(-1)^i(\iota_i)_\#t_{n-1}$.
Łańcuch $z=S(u_n)-u_n-b_n$
jest cyklem: stosując $\partial$,
używamy $\partial S=S\partial$,
równości (*) na ścianach
i $\partial^2=0$.
W wypukłym $\Delta^n$ taki cykl jest
brzegiem: z barycentrum $b$ łączymy
każdy jego singularny sympleks prostym
stożkiem. Dla $\tau:\Delta^k\to\Delta^n$ stożek $b*\tau$
ma wzór
\[
 (b*\tau)(s_0,\ldots,s_{k+1})
 =s_0b+(1-s_0)\tau\left(\frac{s_1}{1-s_0},\ldots,
                         \frac{s_{k+1}}{1-s_0}\right)
 \quad(s_0<1),
\]
z wartością $b$ dla $s_0=1$. Jest ciągły także w tym wierzchołku,
bo obraz $\tau$ jest ograniczony. Brzeg tego stożka to
$\tau-b*(\partial\tau)$ (dla $0$-sympleksu
dodatkowo odejmujemy punkt $b$);
po zsumowaniu nad cyklem składniki
stożkowe się kasują. Daje to $t_n$.
Przenosząc $t_n$ przez dowolne
$\sigma:\Delta^n\to X$, otrzymujemy
operator $T_n(\sigma)=\sigma_\#t_n$ o nośniku zawartym
w $\sigma(\Delta^n)$, który spełnia
\[
 \partial T+T\partial=S-\id.
\]
Zatem $S^r-\id=\partial T^{(r)}+
T^{(r)}\partial$, gdzie
$T^{(r)}=T+ST+\cdots+S^{r-1}T$.

Średnice części po kolejnych podziałach
barycentrycznych dążą do zera: dla
$n$-sympleksu każdy nowy wierzchołek
jest średnią wierzchołków pewnej ściany,
a nowa krawędź łączy barycentra zagnieżdżonych ścian.
Jeśli mają one $p<q\leq n+1$ wierzchołków, to
$b_G=(p/q)b_F+((q-p)/q)b_{G\setminus F}$.
Zatem $\|b_G-b_F\|\leq(q-p)\operatorname{diam}(\Delta)/q
\leq n\operatorname{diam}(\Delta)/(n+1)$.
Średnica wypukłego sympleksu jest największą odległością jego
wierzchołków, więc dla $n\geq1$ otrzymujemy geometryczny spadek
średnic (także na ścianach). Dla $n=0$ sympleks jest już punktem. Dla danego $\sigma$
przeciwobrazy $U_\alpha$ tworzą
otwarte pokrycie zwartego $\Delta^n$;
jego liczba Lebesgue'a jest dodatnia.
Po dostatecznie wielu podziałach
każda część trafia przez $\sigma$
do jednego $U_\alpha$.
Każdy łańcuch jest skończoną sumą,
więc można wybrać wspólne $r$
dla wszystkich jego sympleksów.

Jeśli $z$ jest cyklem, $S^rz$
jest mały i różni się od $z$
o $\partial T^{(r)}z$; to dowodzi
surjektywności na homologii.
Jeśli mały cykl $z$ jest brzegiem
$c$ w pełnym kompleksie, podzielmy
$c$ tak wiele razy, by $S^rc$
był mały. Łańcuch $T^{(r)}z$
też jest mały, ponieważ każda
jego część leży w obrazie
jednego pierwotnego, małego sympleksu.
Równości $\partial c=z$ i
$\partial T^{(r)}z=S^rz-z$ dają
\[
 \partial\bigl(S^rc-T^{(r)}z\bigr)=z.
\]
Zatem $z$ jest brzegiem już
w podkompleksie małych łańcuchów:
otrzymujemy iniektywność.
\end{proof}

Stożkowanie wypełnia każdy cykl dodatniego stopnia w wypukłym
sympleksie, a w stopniu zero tylko cykl o sumie współczynników zero.
W konstrukcji $T$ stopień zero załatwiamy przez $t_0=0$,
bo $S_0=\id$. Lokalność $T$ jest istotna dla dowodu iniektywności.

\begin{twierdzenie}[Wycinanie]
\label{thm:wycinanie}
Jeżeli $Z\subseteq A\subseteq X$ oraz
$\overline Z\subseteq\operatorname{int}_X A$,
to włączenie par
\[
 (X\setminus Z,A\setminus Z)\longrightarrow(X,A)
\]
indukuje izomorfizm na wszystkich
grupach homologii względnej.
\end{twierdzenie}
\begin{proof}
Oznaczmy $U=X\setminus\overline Z$
oraz $V=\operatorname{int}_X A$.
To otwarte pokrycie $X$;
po ograniczeniu do $X\setminus Z$
otrzymujemy pokrycie
$\{U,V\setminus Z\}$.
Lemat~\ref{lem:male-lancuchy} działa
także po wzięciu ilorazu przez
łańcuchy w $A$: izomorfizmy
dla $A$ i $X$, zestawione
z długimi ciągami par, dają
izomorfizm dla odpowiadających
kompleksów względnych
(Lemat~\ref{lem:pieciu}).

Policzmy teraz te małe kompleksy.
Ponieważ $V\subseteq A$,
\[
 C_\bullet^{\{U,V\}}(X)
   =C_\bullet(U)+C_\bullet(V),\qquad
 C_\bullet^{\{U,V\}}(A)
   =C_\bullet(U\cap A)+C_\bullet(V).
\]
Ich iloraz jest kanonicznie
$C_\bullet(U)/C_\bullet(U\cap A)$:
sympleks należący jednocześnie do
$U$ i $V$ należy do $U\cap A$,
więc jądro tej identyfikacji jest
dokładnie wskazanym podkompleksem.
Analogiczny iloraz dla
$(X\setminus Z,A\setminus Z)$
i jego pokrycia jest \emph{tym samym}
kompleksem
$C_\bullet(U)/C_\bullet(U\cap A)$.
Włączenie par staje się na tych
modelach identycznością, co
kończy dowód.
\end{proof}

Wycinanie mówi, że usunięcie części
leżącej głęboko we wnętrzu $A$
nie zmienia tego, co para $(X,A)$
mierzy \emph{poza} $A$. Warunek
o domknięciu ma znaczenie:
zapewnia otwarte pokrycia użyte
w dowodzie.

\begin{twierdzenie}[Ciąg Mayera--Vietorisa]
\label{thm:mv-przestrzenie}
Jeśli $X=U\cup V$, gdzie $U,V$
są otwarte w $X$, to istnieje
długi ciąg dokładny
\begin{equation}\label{eq:mv-przestrzenie}
 \cdots\to H_n(U\cap V)
 \xrightarrow{(i_*,-j_*)}
 H_n(U)\oplus H_n(V)
 \xrightarrow{k_*+\ell_*}H_n(X)
 \xrightarrow{\delta}H_{n-1}(U\cap V)\to\cdots.
\end{equation}
Tak samo działają grupy zredukowane.
\end{twierdzenie}
\begin{proof}
Wewnątrz $C_\bullet(X)$ rozważmy
$B_\bullet=C_\bullet(U)+C_\bullet(V)$.
Mamy krótki ciąg kompleksów
\[
 0\to C_\bullet(U\cap V)
 \xrightarrow{c\mapsto(c,-c)}
 C_\bullet(U)\oplus C_\bullet(V)
 \xrightarrow{(a,b)\mapsto a+b}
 B_\bullet\to0.
\]
Iniektywność pierwszego odwzorowania jest jasna;
jądro drugiej składa się z par
$(c,-c)$, w których każdy singularny
sympleks występujący w $c$ należy
zarazem do $U$ i do $V$.
Surjektywność wynika z definicji sumy.
Ponadto $B_\bullet=C_\bullet^{\{U,V\}}(X)$,
więc Lemat~\ref{lem:male-lancuchy}
identyfikuje $H_*(B)$ z $H_*(X)$.
Zastosowanie Twierdzenia~\ref{thm:les-homologia}
daje~\eqref{eq:mv-przestrzenie}.
W wersji zredukowanej dołączamy
do każdego kompleksu stopień $-1$
równy $\Z$ i różniczkę
$\varepsilon:C_0\to\Z$.
Na stopniu $-1$ ten sam krótki
ciąg ma postać
$0\to\Z\xrightarrow{c\mapsto(c,-c)}
\Z\oplus\Z\xrightarrow{(a,b)\mapsto a+b}
\Z\to0$, nadal dokładną.
W stopniach nieujemnych
obowiązuje poprzedni rachunek.
Homologie tych powiększonych
kompleksów są właśnie
homologiami zredukowanymi;
dotyczy to również pustego
przecięcia, dla którego
$\widetilde H_{-1}(\varnothing)=\Z$.
\end{proof}

Jeśli $c=a+b\in B_n$ jest cyklem,
gdzie $a\in C_n(U)$ i $b\in C_n(V)$,
to $\partial a=-\partial b$ jest
łańcuchem w $U\cap V$, a
$\delta[c]=[\partial a]$ z wyborem
znaku zgodnym z pierwszym odwzorowaniem.
To praktyczna recepta na odwzorowanie łączące.

\begin{przyklad}[Okrąg z dwóch łuków]
\label{prz:mv-okrag}
Pokryjmy $S^1$ dwoma otwartymi,
ściągalnymi łukami $U,V$, których
przecięcie składa się z dwóch
łuków. Mamy
$\widetilde H_0(U)=
\widetilde H_0(V)=0$,
a $\widetilde H_0(U\cap V)\cong\Z$
generuje różnica klas punktów
z dwóch składowych.
Zredukowany ciąg
Mayera--Vietorisa daje
\[
 0=\widetilde H_1(U)\oplus
 \widetilde H_1(V)\to
 \widetilde H_1(S^1)
 \xrightarrow[\cong]{\delta}
 \widetilde H_0(U\cap V)\to0.
\]
Zatem $H_1(S^1)\cong\Z$.
Geometria odwzorowania $\delta$ jest prosta:
rozcinając cykl okręgu na łuki
leżące w $U$ i $V$, otrzymujemy
dwa końce w różnych składowych
ich przecięcia.
\end{przyklad}

\begin{twierdzenie}[Homologia sfer]
\label{thm:homologia-sfer}
Dla $n\geq0$:
\[
 \widetilde H_k(S^n;\Z)=
 \begin{cases}\Z,&k=n,\\0,&k\ne n.\end{cases}
\]
Ponadto dla $n\geq1$
\[
 H_k(D^n,S^{n-1};\Z)=
 \begin{cases}\Z,&k=n,\\0,&k\ne n.\end{cases}
\]
\end{twierdzenie}
\begin{proof}
$S^0$ ma dwie składowe punktowe:
$H_0(S^0)=\Z^2$ oraz
$\widetilde H_0(S^0)=\Z$.
Wyższa homologia znika, bo jest
to suma dwóch kopii homologii
punktu z Przykładu~\ref{prz:homologia-punktu}.

Dla $n\geq1$ weźmy lekko
powiększone otwarte półkule
$U,V\subset S^n$. Każda jest
ściągalna, a $U\cap V$
deformacyjnie zwija się na
sferę równikową $S^{n-1}$.
Zredukowany ciąg
Mayera--Vietorisa, wraz z
zerowaniem homologii $U,V$,
daje dla każdego $k$ izomorfizm
\[
 \widetilde H_k(S^n)
 \xrightarrow{\ \cong\ }
 \widetilde H_{k-1}(S^{n-1}).
\]
W stopniu $k=1$ używa się
właśnie zredukowanego $H_0$;
gdy $n=1$, obie składowe
przecięcia są niezbędne.
Indukcja od $S^0$ daje pierwszy wzór.

Dysk $D^n$ jest ściągalny. W ciągu
zredukowanym pary $(D^n,S^{n-1})$
grupy dysku są zerowe, a odwzorowanie
łączące daje
$H_k(D^n,S^{n-1})\cong
\widetilde H_{k-1}(S^{n-1})$
dla $k\geq1$. Dla $k=0$
odwzorowanie $H_0(S^{n-1})\to H_0(D^n)$
jest surjekcją, więc grupa
względna również znika.
To daje drugi wzór.
\end{proof}

\begin{uwaga}[Klasa fundamentalna i orientacja]
Generator $H_n(D^n,S^{n-1})\cong\Z$
nazywamy \emph{klasą orientacji dysku}.
Wybór jednego z jego dwóch znaków
określa orientację; w długim ciągu
pary $\delta$ posyła ją na
generator $\widetilde H_{n-1}(S^{n-1})$.
Różniczkowe ujęcie orientacji
i całkowania pozostawiamy
osobnemu rozdziałowi.
\end{uwaga}

\section{Od triangulacji do homologii przestrzeni}
\label{sec:porownanie-symplicjalne}

Dla kompleksu symplicjalnego zachowujemy ustalony porządek
wierzchołków, a dla $\Delta$-kompleksu porządki jego map charakterystycznych.
Każdy bazowy sympleks $\sigma:\Delta^n\to|K|$ jest wtedy
jednocześnie singularnym sympleksem. Daje to odwzorowanie łańcuchowe
\[
 I:C_\bullet^\Delta(K)\longrightarrow C_\bullet(|K|),
 \qquad [\sigma]\longmapsto \sigma,
\]
bo oba operatory $\partial$ sumują te same
ściany z tymi samymi znakami.
Mapę $I$ rozszerzamy liniowo: na przykład $I(-[\sigma])=-\sigma$.
Nie zastępujemy prawej strony osobnym generatorem otrzymanym
przez odwrócenie parametryzacji. Nie jest to
izomorfizm grup łańcuchów: po prawej są
wszystkie odwzorowania ciągłe, po lewej tylko
wybrane sympleksy. Niezwykłe jest to,
że po przejściu do homologii różnica znika.

Potrzebny będzie prosty sposób traktowania
komórki z całym brzegiem zwiniętym do punktu.

\begin{lemat}[Iloraz pary komórkowej]
\label{lem:iloraz-cw-homologia}
Niech $A\ne\varnothing$ będzie podkompleksem skończonego
kompleksu CW $X$. Jeśli istnieje otwarte
otoczenie $N$ zbioru $A$, które silnie
deformacyjnie retrahuje się na $A$, to
odwzorowanie ilorazowe indukuje izomorfizm
\[
 H_k(X,A)\xrightarrow{\ \cong\ }
 H_k(X/A,\{*\})=\widetilde H_k(X/A)
 \quad(k\geq0).
\]
Takie otoczenie istnieje, gdy
$A=X^{n-1}\ne\varnothing$ i $X=X^n$ jest skończonym
szkieletem CW dla $n\geq1$, a także gdy do skończonego $A$
dołączamy pojedynczą komórkę dodatniego wymiaru.
\end{lemat}
\begin{proof}
Włączenie $A\hookrightarrow N$ jest
równoważnością homotopijną. Porównując
długie ciągi par $(X,A)$ i $(X,N)$
za pomocą Lematu~\ref{lem:pieciu},
dostajemy $H_k(X,A)\cong H_k(X,N)$.
Po zwinięciu $A$ do punktu homotopia schodzi na $N/A$.
Uzasadnijmy jej ciągłość również względem czasu.
Skończony kompleks CW jest zwarty, bo jest obrazem skończonej
sumy zwartych dysków; domknięty $A$ też jest zwarty.
Projekcja $q:N\to N/A$ jest domknięta: dla domkniętego $F\subseteq N$
zbiór $q^{-1}q(F)$ jest równy $F$ lub $F\cup A$.
Ma zwarte włókna, więc $q\times\id_I$ jest również domknięta.
Istotnie, jeśli domknięty zbiór omija włókno nad $(y,t)$,
skończone pokrycie tego włókna prostokątami otwartymi daje
otoczenia $O\supseteq q^{-1}(y)$ i $J\ni t$ z $O\times J$
rozłącznym z tym zbiorem; domkniętość $q$ daje otoczenie
$W\ni y$ takie, że $q^{-1}(W)\subseteq O$.
Stąd $(q\times\id_I)$ jest odwzorowaniem ilorazowym,
co pozwala przenieść homotopię. Zatem $N/A$ silnie zwija się do $\{*\}$,
więc tak samo
$H_k(X/A,\{*\})\cong H_k(X/A,N/A)$.

Porównajmy środkowe grupy.
Zbiór $A$ jest domknięty
(dla skończonego CW wynika to
z topologii słabej), zatem
$\{X\setminus A,N\}$ jest
otwartym pokryciem $X$.
W ilorazie mamy pokrycie
$\{(X/A)\setminus\{*\},N/A\}$.
Lemat o małych łańcuchach,
zastosowany również do $N$,
pozwala obliczać homologie
tych dwóch par odpowiednio za pomocą kompleksów
\[
 \frac{C_\bullet(X\setminus A)}
      {C_\bullet(N\setminus A)}
 \quad\text{oraz}\quad
 \frac{C_\bullet((X/A)\setminus\{*\})}
      {C_\bullet((N/A)\setminus\{*\})}.
\]
Poza $A$ odwzorowanie ilorazowe jest
homeomorfizmem, więc te dwa
kompleksy są izomorficzne.
Składając trzy izomorfizmy,
otrzymujemy tezę.

Dla $A=X^{n-1}\subset X^n$
weźmy $N$ złożony z $A$ oraz,
w każdej dołączonej $n$-komórce,
obrazu pierścienia
$\{x\in D^n:\|x\|>1/2\}$.
Jest otwarty w topologii
nadanej przez odwzorowania dołączające.
W każdym pierścieniu promieniowo
przesuwamy $r$ od wartości
$r>1/2$ do $1$, pozostawiając
brzeg i cały $A$ nieruchome.
Na częściach dyskowych używamy wzoru
$ru\mapsto((1-t)r+t)u$. Ciągłość przy $A$ sprawdzamy jak
w dowodzie Twierdzenia~\ref{thm:pi1-dolaczenie-komorki}:
w każdym otoczeniu zachowujemy tylko te punkty pasa, których
cały odcinek radialny do brzegu mieści się w tym otoczeniu.
Te mniejsze zbiory są otwarte i pozostają w pierwotnym otoczeniu
przez cały czas homotopii. Daje to silną retrakcję $N\to A$.
Ten sam argument działa przy dołączeniu pojedynczej komórki.
\end{proof}

\begin{lemat}[Generator pojedynczego sympleksu]
\label{lem:orientacja-simpleksu}
Identycznościowe odwzorowanie singularne
$u_n:\Delta^n\to\Delta^n$ określa
generator
$H_n(\Delta^n,\partial\Delta^n)\cong\Z$
dla $n\geq1$.
\end{lemat}
\begin{proof}
Jest cyklem względnym, bo jego
brzeg leży w $\partial\Delta^n$.
Brzeg sympleksu jest sferą
(Stwierdzenie~\ref{prop:brzeg-sympleksu}),
a jej homologię policzyliśmy
w Twierdzeniu~\ref{thm:homologia-sfer}.
Wynika stąd, że odwzorowanie łączące
\[
 \delta:H_n(\Delta^n,\partial\Delta^n)
 \longrightarrow
 \widetilde H_{n-1}(\partial\Delta^n)
 \]
jest izomorfizmem i posyła
$[u_n]$ na klasę cyklu
$\partial u_n$.
Pokażemy indukcyjnie, że ten
cykl jest generatorem.
Dla $n=1$ jest to różnica
końców, generator
$\widetilde H_0(S^0)$.
W wyższym wymiarze rozdzielmy
$\partial\Delta^n$ na jedną
ścianę $F\cong\Delta^{n-1}$
i sumę $B$ pozostałych ścian.
$F$ i $B$ są ściągalne:
$B$ jest stożkiem z przeciwległego
wierzchołka nad $\partial F$;
ich przecięciem jest $\partial F$.
W ciągu Mayera--Vietorisa
odwzorowanie łączące
$\widetilde H_{n-1}(\partial\Delta^n)
\to\widetilde H_{n-2}(\partial F)$
jest izomorfizmem. Aby zastosować wersję dla zbiorów otwartych, przedstawmy
oba dyski $F,B$ jako stożki nad ich wspólnym brzegiem $\partial F$.
Niech $r\in[0,1]$ oznacza w każdym z nich parametr radialny,
z $r=1$ na $\partial F$. Zbiory
$U=F\cup\{x\in B:r(x)>1/2\}$ oraz
$V=B\cup\{x\in F:r(x)>1/2\}$ są otwarte w $F\cup B$:
na każdej z dwóch części ich przeciwobrazy są otwarte.
Przesuwając dodatkowe pasy do $r=1$, retrahujemy je na $F,B$,
a ich przecięcie na $\partial F$.
Cykl $\partial u_n$ jest sumą
$\pm u_F$ na $F$ i pozostałych
ścian na $B$; jego obrazem przez
odwzorowanie łączące jest
$\pm[\partial u_F]$.
To generator na mocy
założenia indukcyjnego.
Zatem $\partial u_n$, a więc
również $u_n$ względnie,
jest generatorem.
\end{proof}

\begin{twierdzenie}[Porównanie homologii]
\label{thm:symplicjalna-singularna}
Dla skończonego kompleksu symplicjalnego
$K$ naturalne odwzorowanie indukuje
\[
 I_*:H_n^\Delta(K)\xrightarrow{\ \cong\ }
 H_n(|K|;\Z)
 \qquad(n\geq0).
\]
To samo zachodzi dla skończonego
$\Delta$-kompleksu, jeśli generatory
po lewej to jego sympleksy
charakterystyczne z ich ścianami.
\end{twierdzenie}
\begin{proof}
Uporządkujmy sympleksy tak, aby
każdy pojawiał się po wszystkich
swoich ścianach. Budujemy kompleks
po jednym sympleksie:
$L\subset K=L\cup_\varphi\Delta^n$.
Załóżmy indukcyjnie, że porównanie
jest izomorfizmem dla $L$.
Iloraz kompleksów symplicjalnych
$C^\Delta(K)/C^\Delta(L)$
ma jeden generator w stopniu $n$
i zero w pozostałych stopniach,
więc jego homologia jest $\Z$
tylko w stopniu $n$.

Jeśli $n\geq1$, po zwinięciu $L$
do punktu przestrzeń $|K|/|L|$
jest $\Delta^n/\partial\Delta^n
\cong S^n$; analogiczny opis
działa dla $\Delta$-kompleksu,
nawet gdy jego ściany były
między sobą utożsamione.
Lemat~\ref{lem:iloraz-cw-homologia}
i Twierdzenie~\ref{thm:homologia-sfer}
dają $H_q(|K|,|L|)=\Z$ dla
$q=n$ i zero w innych stopniach.
Odwzorowanie względne przenosi jedyny
generator na klasę charakterystycznego sympleksu.
Mapa par $(\Delta^n,\partial\Delta^n)\to(|K|,|L|)$
indukuje homeomorfizm ich ilorazów: jest ciągłą bijekcją
ze zwartej przestrzeni do przestrzeni Hausdorffa.
Naturalność map do ilorazów w Lemacie~\ref{lem:iloraz-cw-homologia}
pokazuje, że także na homologii względnej jest izomorfizmem.
Na mocy Lematu~\ref{lem:orientacja-simpleksu} przenosi więc
generator na generator, a nie na jego dowolną wielokrotność.
Dla $n=0$ nowy wierzchołek jest
osobną składową i odpowiednie
odwzorowanie względne jest izomorfizmem
$\Z\to\Z$ w stopniu $0$.

Mamy teraz przemienny diagram
długich ciągów par dla kompleksów
symplicjalnych i singularnych.
Na wyrazach dotyczących $L$
odwzorowania pionowe są izomorfizmami
z założenia indukcyjnego,
a na względnych właśnie
to pokazaliśmy. Lemat~\ref{lem:pieciu}
zastosowany do kolejnych pięciu
wyrazów daje izomorfizm dla $K$.
Zaczynamy od kompleksu pustego,
dla którego wszystkie zwykłe
grupy homologii są zerowe.
\end{proof}

W praktyce wolno zatem obliczać
homologię przestrzeni z jednej
wygodnej triangulacji: otrzymany
wynik nie zależy od wyboru
triangulacji. Podziały sympleksów
zmieniają grupy łańcuchów,
ale nie grupy homologii.

\section{Homologia komórkowa}
\label{sec:homologia-komorkowa}

W triangulacji torusa potrzeba
trójkątów i ich ścian. Struktura
CW z Sekcji~\ref{sec:cw-kompleksy}
może opisywać go jednym punktem,
dwoma pętlami i jedną
$2$-komórką. Homologia komórkowa
zachowuje tę oszczędność.

\begin{twierdzenie}[Względna homologia warstwy komórek]
\label{thm:warstwa-komorek}
Niech $X$ będzie skończonym
kompleksem CW, $X^n$ jego
$n$-szkieletem i $X^{-1}=\varnothing$.
Dla $n\geq1$:
\[
 H_k(X^n,X^{n-1})
 \cong
 \begin{cases}
 \displaystyle\bigoplus_{\alpha\in I_n}\Z[e_\alpha^n],
 &k=n,\\[3pt]0,&k\ne n,
 \end{cases}
\]
gdzie $I_n$ oznacza zbiór $n$-komórek.
Dla $n=0$ mamy
$H_0(X^0)=\bigoplus_{\alpha\in I_0}\Z[e_\alpha^0]$
i $H_k(X^0)=0$ dla $k\ne0$.
\end{twierdzenie}
\begin{proof}
Jeśli $X^{n-1}=\varnothing$ i $n\geq1$, to również $X^n=\varnothing$:
nie można dołączyć dysku o niepustym brzegu do pustej przestrzeni.
Teza jest wtedy oczywista. W pozostałych przypadkach
w ilorazie $X^n/X^{n-1}$
brzeg każdej dołączonej kopii
$D^n$ staje się jednym punktem,
a wnętrza komórek pozostają
rozłączne. Z definicji topologii
CW otrzymujemy bukiet
$\bigvee_{\alpha\in I_n}
(D^n/\partial D^n)\cong
\bigvee_{\alpha\in I_n}S^n$.
Lemat~\ref{lem:iloraz-cw-homologia}
zamienia homologię względną
na zredukowaną homologię bukietu.
Dla skończonego bukietu
zredukowana homologia jest sumą
zredukowanych homologii składników:
dołączamy kolejne sfery. Dla bukietu $Y\vee S^n$
bierzemy $U=Y\cup B$ i $V=S^n\cup N$, gdzie $B$ jest małą
otwartą kulą wokół punktu wspólnego w nowej sferze, a $N$
jest sumą takich kul w dotychczasowych sferach.
Te zbiory są otwarte w topologii ilorazowej bukietu,
retrahują się na $Y,S^n$, a $U\cap V=N\vee B$ jest ściągalne.
Ciąg Mayera--Vietorisa daje izomorfizm indukowany włączeniami
składników, a więc właśnie ich sumę prostą.
Twierdzenie~\ref{thm:homologia-sfer}
daje wskazane grupy.
Znaki generatorów zależą
od wyboru orientacji komórek.
Dla $n=0$ mamy skończony
zbiór punktów; rachunek wynika
z homologii punktu.
\end{proof}

\begin{definicja}[Kompleks komórkowy]
\index{kompleks komorkowy@Kompleks komórkowy}
\label{def:kompleks-komorkowy}
Kładziemy
$C_n^{\mathrm{CW}}(X)=H_n(X^n,X^{n-1})$.
Dla $n\geq1$ jego operator brzegu
jest złożeniem
\begin{equation}\label{eq:cellular-boundary}
 H_n(X^n,X^{n-1})
 \xrightarrow{\ \delta_n\ }
 H_{n-1}(X^{n-1})
 \xrightarrow{\ j_{n-1}\ }
 H_{n-1}(X^{n-1},X^{n-2}),
\end{equation}
gdzie $\delta_n$ pochodzi z
długiego ciągu pary, a
$j_{n-1}$ z ilorazu łańcuchów.
Przyjmujemy $\partial_0^{\mathrm{CW}}=0$.
\end{definicja}

\begin{twierdzenie}[Komórki obliczają homologię]
\label{thm:komorkowa-singularna}
Operator~\eqref{eq:cellular-boundary}
spełnia $(\partial^{\mathrm{CW}})^2=0$.
Dla skończonego CW istnieją
izomorfizmy
\[
 H_n(C_\bullet^{\mathrm{CW}}(X),
       \partial^{\mathrm{CW}})
 \cong H_n(X;\Z)
 \qquad\text{dla każdego }n\geq0.
\]
\end{twierdzenie}
\begin{proof}
Dla $n=1$ równość brzegu brzegu wynika z $\partial_0^{\mathrm{CW}}=0$.
Dla $n\geq2$ z dokładności długiego ciągu
pary $(X^{n-1},X^{n-2})$
wynika $\delta_{n-1}j_{n-1}=0$.
Stąd
\[
 \partial^{\mathrm{CW}}_{n-1}
 \partial^{\mathrm{CW}}_n
 =j_{n-2}\delta_{n-1}j_{n-1}\delta_n=0.
\]

Rozpiszmy porównanie homologii.
Indukcja po szkieletach,
z użyciem długiego ciągu pary
i Twierdzenia~\ref{thm:warstwa-komorek},
daje $H_k(X^m)=0$ dla $k>m$.
Włączenie $X^{m-1}\hookrightarrow X^m$
indukuje ponadto izomorfizm na
$H_k$ dla $k<m-1$: dwa sąsiednie
wyrazy względne o stopniach
$k+1$ i $k$ są wtedy zerowe.
W szczególności komórki wymiaru
$\geq n+2$ nie zmieniają $H_n$.

Dla $n\geq1$ w ciągu pary
$(X^n,X^{n-1})$ mamy odcinek
\[
 0=H_n(X^{n-1})\longrightarrow
 H_n(X^n)\xrightarrow{\ j_n\ }
 C_n^{\mathrm{CW}}
 \xrightarrow{\ \delta_n\ }
 H_{n-1}(X^{n-1}).
\]
Odwzorowanie $j_{n-1}:H_{n-1}(X^{n-1})
\to C_{n-1}^{\mathrm{CW}}$
jest iniektywne: jego jądro
jest obrazem
$H_{n-1}(X^{n-2})=0$.
Dla $n=1$ używamy
$X^{-1}=\varnothing$, co daje
ten sam wniosek.
Zatem $\ker\partial_n^{\mathrm{CW}}
=\ker\delta_n=\operatorname{im}j_n$
i $j_n$ identyfikuje cykle
komórkowe z $H_n(X^n)$.

W następnym długim ciągu,
dla $(X^{n+1},X^n)$,
\[
 C_{n+1}^{\mathrm{CW}}
 \xrightarrow{\ \delta_{n+1}\ }
 H_n(X^n)\longrightarrow
 H_n(X^{n+1})\longrightarrow
 H_n(X^{n+1},X^n)=0.
\]
Jądro przejścia do $H_n(X^{n+1})$
jest więc $\operatorname{im}\delta_{n+1}$.
Po włożeniu w $C_n^{\mathrm{CW}}$
przez $j_n$ jest to dokładnie
$\operatorname{im}\partial_{n+1}^{\mathrm{CW}}$.
Otrzymujemy
$H_n^{\mathrm{CW}}(X)\cong
H_n(X^{n+1})\cong H_n(X)$.
Dla $n=0$ ten sam argument działa
z $H_0(X^0)=C_0^{\mathrm{CW}}$:
$H_0(X^1)$ jest ilorazem przez
obrazy brzegów $1$-komórek,
a wyższe komórki już nie
zmieniają $H_0$.
\end{proof}

\section{Odwzorowania dołączające i konkretne obliczenia}
\label{sec:obliczenia-homologii}

\begin{definicja}[Stopień odwzorowania sfer]
\index{stopien odwzorowania sfer@Stopień odwzorowania sfer}
\label{def:stopien-sfer}
Dla $m\geq1$ każde ciągłe odwzorowanie
$h:S^m\to S^m$ indukuje
endomorfizm $H_m(S^m)\cong\Z$.
Istnieje więc dokładnie jedna
liczba całkowita $\deg h$ taka,
że $h_*([S^m])=(\deg h)[S^m]$.
Nazywamy ją \emph{stopniem} $h$.
\end{definicja}

Z definicji i funktorialności
$\deg(\id)=1$ oraz
$\deg(h\circ g)=\deg h\deg g$.
Homotopijne odwzorowania mają ten sam
stopień. Na $S^1=\{z\in\C:|z|=1\}$
odwzorowanie $z\mapsto z^r$ ma stopień $r\in\Z$.
Uzasadnijmy przejście od liczby obiegów do homologii.
Pętla $\gamma(t)=e^{2\pi it}$ reprezentuje generator $H_1(S^1)$:
po podziale na łuki jej obraz przez izomorfizm Mayera--Vietorisa
z Przykładu~\ref{prz:mv-okrag} jest różnicą punktów dwóch składowych przecięcia.
Dla dowolnych dwóch pętli $\alpha,\beta$ opartych w tym samym punkcie
mamy $[\alpha*\beta]=[\alpha]+[\beta]$ w $H_1$.
Istotnie, złożenie $\alpha*\beta$ z afinicznym odwzorowaniem
$\Delta^2\to[0,1]$ o wartościach $0,1/2,1$ w kolejnych wierzchołkach
jest singularnym $2$-sympleksem o brzegu
$\beta-(\alpha*\beta)+\alpha$.
Ponadto homotopia pętli względem końców daje homologiczne cykle
przez wzór pryzmatyczny; pętla stała jest brzegiem stałego $2$-sympleksu.
Stąd $[\bar\alpha]=-[\alpha]$, bo $\alpha*\bar\alpha$ jest
homotopijna do stałej. Pętla $t\mapsto e^{2\pi irt}$ jest,
z dokładnością do zmiany parametryzacji, złożeniem $r$ obiegów dla $r>0$,
$|r|$ obiegów przeciwnych dla $r<0$
i jest stała dla $r=0$. Jej klasa to zatem $r[\gamma]$.

\begin{twierdzenie}[Współczynniki brzegu komórkowego]
\label{thm:formula-komorkowa}
Niech $\varphi_\alpha:S^{n-1}\to X^{n-1}$
będzie odwzorowaniem dołączającym
zorientowanej $n$-komórki
$e_\alpha^n$, gdzie $n\geq2$.
Dla każdej zorientowanej
$(n-1)$-komórki $e_\beta^{n-1}$
zwińmy w $X^{n-1}$ do punktu
dopełnienie jej otwartego wnętrza.
Otrzymamy odwzorowanie
\[
 q_\beta:X^{n-1}\longrightarrow
 X^{n-1}/(X^{n-1}\setminus e_\beta^{n-1})
 \cong S^{n-1}.
\]
Sferę docelową orientujemy zgodnie z komórką $e_\beta^{n-1}$,
a dziedzinę $\varphi_\alpha$ jako brzeg zorientowanej komórki $e_\alpha^n$.
Wówczas
\begin{equation}\label{eq:formula-komorkowa}
 \partial_n^{\mathrm{CW}}(e_\alpha^n)
 =\sum_{\beta}
 \deg(q_\beta\circ\varphi_\alpha)
 e_\beta^{n-1}.
\end{equation}
W stopniu $1$ brzeg zorientowanej
komórki od początku $v_-$ do końca
$v_+$ wynosi $\partial_1e=v_+-v_-$.
\end{twierdzenie}
\begin{proof}
Generator $e_\alpha^n$ jest
obrazem klasy orientacji
$(D^n,S^{n-1})$ w parze
$(X^n,X^{n-1})$. Pierwsza
strzałka~\eqref{eq:cellular-boundary}
zwraca klasę obrazu
zorientowanego brzegu $S^{n-1}$
przez $\varphi_\alpha$:
tak działa odwzorowanie łączące,
bo na łańcuchach bierze
dosłowny brzeg podniesionego
cyklu względnego.
Druga strzałka przenosi tę
klasę do
$H_{n-1}(X^{n-1},X^{n-2})$.
Według
Twierdzenia~\ref{thm:warstwa-komorek}
jest to suma kopii $\Z$,
po jednej dla każdego $\beta$.
Rzutowanie na składnik $\beta$
jest, poprzez iloraz z
Lematu~\ref{lem:iloraz-cw-homologia},
tym samym co $q_{\beta*}$.
Zatem współczynnik przy
$e_\beta^{n-1}$ to liczba,
przez którą
$(q_\beta\circ\varphi_\alpha)_*$
mnoży generator homologii
$S^{n-1}$: z definicji
jest to jego stopień.
Dla $n=1$ używamy
bezpośrednio
$\partial[0,1]=[1]-[0]$;
oba końce mogą w $X^0$
oznaczać ten sam wierzchołek.
\end{proof}

Warto przeczytać wzór geometrycznie:
liczymy, ile razy, z jakimi znakami,
brzeg wyższej komórki przechodzi
przez daną komórkę niższą. Komórka
z niezerowym współczynnikiem brzegu
może zabić klasę homologii.

\begin{figure}[!htbp]
\centering
\begin{tikzpicture}[>=Stealth,line cap=round,line join=round]
 \node[draw=manifoldblue!65!black,fill=manifoldblue!9,
  rounded corners=4pt,minimum width=2.15cm,
  minimum height=.8cm] (c2) at (1.2,0) {$C_2=\Z$};
 \node[draw=manifoldgreen!65!black,fill=manifoldgreen!9,
  rounded corners=4pt,minimum width=2.35cm,
  minimum height=.8cm] (c1) at (4.5,0) {$C_1=\Z^2$};
 \node[draw=manifoldgold!75!black,fill=manifoldgold!10,
  rounded corners=4pt,minimum width=2.15cm,
  minimum height=.8cm] (c0) at (7.8,0) {$C_0=\Z$};
 \draw[->,thick,manifoldred] (c2)--(c1)
  node[midway,above=5pt] {$\partial_2=0$};
 \draw[->,thick,manifoldred] (c1)--(c0)
  node[midway,above=5pt] {$\partial_1=0$};
 \node[align=center,font=\small,manifoldblue!75!black]
  at (1.2,-.9) {jedna\\$2$-komórka};
 \node[align=center,font=\small,manifoldgreen!75!black]
  at (4.5,-.9) {dwie\\pętle};
 \node[font=\small,manifoldgold!75!black]
  at (7.8,-.9) {punkt};
\end{tikzpicture}
\caption{Kompleks komórkowy torusa. Ma trzy
niezerowe grupy łańcuchów, a oba odwzorowania
brzegu znikają: pętle kończą się w tym
samym punkcie, a słowo dołączające
komórkę $2$-wymiarową jest komutatorem.}
\label{fig:kompleks-torusa}
\end{figure}

\begin{przyklad}[Sfera z dwóch komórek]
\label{prz:homologia-cw-sfera}
W modelu z
Przykładu~\ref{prz:cw-s2}
jest jedna komórka $e^0=v$
i jedna $e^2$, bez $1$-komórek.
Kompleks komórkowy to
\[
 0\longrightarrow\Z e^2
 \xrightarrow{\;0\;}0
 \xrightarrow{\;0\;}\Z v
 \longrightarrow0.
\]
Stąd $H_2(S^2)=\Z$, $H_1(S^2)=0$,
$H_0(S^2)=\Z$. Ogólniej
$S^n=D^n/\partial D^n$ ma
jedną komórkę w stopniu $0$
i jedną w stopniu $n$
($n\geq2$). Dla $S^1$
jedna $1$-komórka jest pętlą,
więc $\partial_1=0$.
W ten sposób odzyskujemy
Twierdzenie~\ref{thm:homologia-sfer}.
\end{przyklad}

\begin{przyklad}[Torus]
\label{prz:homologia-cw-torus}
Model CW z
Przykładu~\ref{prz:pi1-torus-rp2}
ma komórki $v$,
$a,b$ oraz $e^2$.
Krawędzie są pętlami, więc
$\partial_1a=\partial_1b=0$.
Brzeg $e^2$ jest dołączony
według słowa $aba^{-1}b^{-1}$.
Współczynnik przy $a$
to $1-1=0$, a przy $b$
to $1-1=0$; jest to
szczególny przypadek
Wzoru~\eqref{eq:formula-komorkowa}.
Zatem $\partial_2=0$ i
\[
 H_0(T^2)=\Z,\qquad
 H_1(T^2)=\Z^2,\qquad
 H_2(T^2)=\Z,\qquad
 H_k(T^2)=0\quad(k>2).
\]
Rysunek~\ref{fig:kompleks-torusa}
pokazuje cały rachunek.
Grupa $\Z^2$ rozpoznaje
dwa niezależne kierunki pętli.
\end{przyklad}

\begin{przyklad}[Płaszczyzna rzutowa i torsja]
\label{prz:homologia-rp2}
Przestrzeń $\mathbb{RP}^2$
ma po jednej komórce
$e^0,e^1,e^2$.
Jedyna $1$-komórka jest pętlą,
więc $\partial_1=0$.
Odwzorowanie dołączające $e^2$
okrąża $\mathbb{RP}^1\cong S^1$
dwa razy, co było widoczne
w Przykładzie~\ref{prz:pi1-torus-rp2}.
Ma stopień $2$ po ustaleniu
orientacji, zatem
\[
 0\longrightarrow\Z
 \xrightarrow{\;\times2\;}\Z
 \xrightarrow{\;0\;}\Z
 \longrightarrow0.
\]
Jądro mnożenia przez $2$
jest zerowe, a iloraz
$\Z/2\Z$ ma element
niezerowy, który po
dwukrotnym dodaniu znika.
Otrzymujemy
\[
 H_0(\mathbb{RP}^2)=\Z,\quad
 H_1(\mathbb{RP}^2)=\Z/2\Z,\quad
 H_2(\mathbb{RP}^2)=0.
\]
To konkretny przykład \emph{torsji}
zdefiniowanej w
Sekcji~\ref{sec:wspolczynniki-euler}.
\end{przyklad}

\begin{uwaga}[Homologia ze współczynnikami]
\label{rem:homologia-wspolczynniki-przestrzeni}
Uzasadnijmy zmianę współczynników także w dodatniej charakterystyce.
Dla przemiennego pierścienia $R$ z jedynką definiujemy $C_n(X;R)$
jako wolny $R$-moduł na singularnych sympleksach, z tym samym wzorem brzegu.
Równoważnie jest to $C_n(X;\Z)\otimes_{\Z}R$.
Homotopie pryzmatyczne i podział barycentryczny mają całkowite współczynniki,
więc wszystkie ich tożsamości pozostają prawdziwe nad $R$.
Ciągi łańcuchów par oraz małych łańcuchów Mayera--Vietorisa są
rozszczepione w każdym stopniu na wybranych bazach, więc są dokładne także nad $R$.
Rachunek punktu, sfer i kolejnych szkieletów przebiega wtedy tak samo:
$C_n^{\mathrm{CW}}(X;R)$ ma po jednej kopii $R$ na komórkę.
Naturalność odwzorowań łączących względem $\Z\to R$ pokazuje,
że całkowity współczynnik brzegu $d$ staje się $d\cdot1_R$.
To uzasadnia obliczenia komórkowe nad $R$; nie zakładamy,
że utworzenie iloczynu tensorowego dowolnego ciągu dokładnego
z $R$ zawsze zachowuje dokładność.
\end{uwaga}

\begin{przyklad}[Te same przestrzenie nad różnymi ciałami]
\label{prz:zmiana-wspolczynnikow-rp2-torus}
Wstawmy teraz do powyższych kompleksów komórkowych
w miejsce $\Z$ dowolne
$\mathbb F\in\{\Q,\R,\C\}$.
Dla płaszczyzny rzutowej dostajemy
\[
 0\longrightarrow\mathbb F
 \xrightarrow{\;\times2\;}\mathbb F
 \xrightarrow{\;0\;}\mathbb F
 \longrightarrow0.
\]
Ponieważ $2$ ma w $\mathbb F$ odwrotność
$1/2$, środkowe odwzorowanie jest
izomorfizmem. Zatem
\[
 H_0(\mathbb{RP}^2;\mathbb F)=\mathbb F,
 \qquad H_k(\mathbb{RP}^2;\mathbb F)=0\quad(k>0).
\]
Klasa $\Z/2$ zniknęła. Dla torusa oba
operatory brzegu nadal są zerowe, więc
\[
 H_0(T^2;\mathbb F)=\mathbb F,\qquad
 H_1(T^2;\mathbb F)=\mathbb F^2,\qquad
 H_2(T^2;\mathbb F)=\mathbb F.
\]
Wymiary tych przestrzeni, odpowiednio
$1,2,1$, to liczby Betti'ego torusa.
Tak działa ogólny
Wzór~\eqref{eq:homologia-char-zero}.
Dla porównania nad $\Z/2$ strzałka
$\times2$ w kompleksie $\mathbb{RP}^2$
staje się zerowa. Wtedy oprócz $H_0$
otrzymujemy również
$H_1(\mathbb{RP}^2;\Z/2)\cong\Z/2$
oraz $H_2(\mathbb{RP}^2;\Z/2)\cong\Z/2$.
\end{przyklad}

\begin{przyklad}[Bukiet okręgów i charakterystyka Eulera]
\label{prz:homologia-bukiet}
Bukiet $r$ okręgów ma jeden
wierzchołek i $r$ pętli,
bez komórek wyższych:
$C_1^{\mathrm{CW}}=\Z^r$,
$C_0^{\mathrm{CW}}=\Z$ i
$\partial_1=0$.
Stąd $H_1=\Z^r$, $H_0=\Z$,
a wyższe homologie znikają.
Dla skończonego CW z
$c_n$ komórkami w wymiarze $n$
równanie Eulera z
Twierdzenia~\ref{thm:euler-homologia}
przyjmuje postać
\[
 \sum_n(-1)^n c_n
 =\sum_n(-1)^n
   \operatorname{rank}_{\Z} H_n(X;\Z).
\]
W szczególności dla torusa
$1-2+1=0$, a dla
$\mathbb{RP}^2$ mamy
$1-1+1=1$; torsja
$\Z/2\Z$ ma rangę zero.
\end{przyklad}

\section{Lokalna homologia i dalsza droga}
\label{sec:lokalna-homologia}

\begin{stwierdzenie}[Lokalna homologia rozmaitości]
\label{prop:lokalna-homologia}
Jeśli $M$ jest topologiczną
rozmaitością wymiaru $m$
bez brzegu i $p\in M$, to
\[
 H_k(M,M\setminus\{p\};\Z)
 \cong
 \begin{cases}\Z,&k=m,\\0,&k\ne m.\end{cases}
\]
Wybór generatora grupy w stopniu $m$
jest \emph{lokalną orientacją}
w punkcie $p$.
\end{stwierdzenie}
\begin{proof}
Wybierzmy otwarte otoczenie $U\subseteq M$ i mapę
$\psi:U\to B^m$ na otwartą kulę jednostkową,
z $\psi(p)=0$. Zbiór
$Z=M\setminus U$
jest domknięty i leży
we wnętrzu $M\setminus\{p\}$,
więc wycinanie
(Twierdzenie~\ref{thm:wycinanie})
daje
\[
 H_k(M,M\setminus\{p\})
 \cong H_k(U,U\setminus\{p\})
\cong H_k(B^m,B^m\setminus\{0\}).
\]
Dla $m\geq1$ kula jest
ściągalna, a kula pozbawiona
środka silnie deformacyjnie retrahuje się na sferę
o promieniu $1/2$, homeomorficzną z $S^{m-1}$:
używamy $ru\mapsto((1-t)r+t/2)u$, co pozostaje w $0<r<1$.
Długi ciąg pary i
Twierdzenie~\ref{thm:homologia-sfer}
dają wynik. Dla $m=0$
punkt jest otwartą składową
mapy lokalnej, a grupa
$H_0(\{p\},\varnothing)=\Z$.
\end{proof}

Istnienie lokalnego generatora
nie oznacza jeszcze, że można
wybrać generatory \emph{spójnie}
na całej rozmaitości. To będzie
treść pojęcia orientowalności.
Na obecnym etapie mamy
już narzędzia do dalszej nauki:
kompleksy łańcuchowe i długie ciągi,
homotopijną niezmienniczość,
wycinanie, Mayera--Vietorisa
oraz model komórkowy, w którym
odwzorowania dołączające kontrolują
operator brzegu. Pozwalają one
później opisać, jak dołączenie
uchwytu zmienia homologię,
bez wprowadzania tutaj
teorii Morse'a i chirurgii.

\par\smallskip
{\footnotesize\noindent\emph{Dalsza lektura:}
Pełne rozwinięcie homologii symplicjalnej,
singularnej i komórkowej, wraz ze
szczegółami technicznymi dotyczącymi
podziału barycentrycznego, znajduje się
w Allen Hatcher,
\emph{Algebraic Topology}, rozdział 2,
\url{https://pi.math.cornell.edu/~hatcher/AT/AT.pdf}.\par}
\clearpage
